\documentclass[11pt]{article}
\usepackage[T1]{fontenc}
\usepackage{lmodern}
\usepackage[margin=2.5cm]{geometry}
\usepackage{booktabs}
\usepackage{amsmath}
\usepackage[hidelinks]{hyperref}
\usepackage{xcolor}

\newcommand{\code}[1]{\texttt{#1}}
\emergencystretch=3em
\newcommand{\wk}{WUKONG}

\title{UDA-Bench NBA experiments\\\large E1 graph structure, X6 coupling ablation, schema-based extractors}
\author{}
\date{October 9, 2026}

\begin{document}
\maketitle

\section{Goal}

Three studies on the same corpus, with the same language model:
\begin{itemize}
  \item \textbf{E1, graph structure.} \wk{}, LightRAG and GraphRAG build a graph from the corpus. We
    measure how the graphs can be navigated and queried, with statistics that can be read without
    knowing the domain: size, node and relation vocabularies, edge direction and properties. Recall
    against the gold tables is a control, to check that the graphs contain comparable information.
  \item \textbf{X6, coupling ablation.} \wk{} with six schemas, from the schema core alone (types,
    fields, endpoints) to the full extraction schema, adding one family of annotations at a time.
  \item \textbf{Schema-based extractors.} Three tools that take a schema as input (Neo4j's
    \code{neo4j-graphrag}, LlamaIndex's \code{SchemaLLMPathExtractor}, OntoGPT's SPIRES), given our
    schema in their own formats, scored like the X6 variants.
\end{itemize}

\section{Common setup}

\paragraph{Corpus and gold.} The NBA dataset of UDA-Bench: 216 English Wikipedia articles,
141 about players, 30 about teams, 29 about cities and 16 about team owners. The gold tables
(player, team, city, manager) have one row per article; the row ID is the document number.

\paragraph{Model.} Every system uses \code{gpt-5.6-luna} with low reasoning effort, and
\code{text-\allowbreak embedding-\allowbreak 3-small} where it embeds. Output tokens below include reasoning tokens.

\paragraph{\wk.} Engine at commit \code{baeeb13}, chunks of 800 tokens with 120 overlap, document
level reading the first 8000 tokens of each document, export to Neo4j CSV
(\code{x6-config.toml}). The full schema (\code{workspace-x6-s4/}) has four entity types (Player,
Team, City, Owner) with typed fields and three relationship types (PlaysFor, with a
\code{tenure} field; LocatedIn; Owns, with a \code{since\_year} field). Each \wk{} configuration
was run three times; tables give the mean and, where runs differ, the range.

\paragraph{Schema hygiene.} The first version of the schema (\code{workspace/}) used example values
taken from the corpus, including gold answers (two players' birth dates, two cities' states, a
team's owner, and subject names), and one instruction used a subject team as its format example.
We replaced every such value by one that is neither a gold cell nor the name of a subject entity,
leaving everything else unchanged. \code{make\_schema\_variants.py} and
\code{make\_coupling\_variants.py} reject any schema whose examples, descriptions or instructions
contain a gold value or a subject's name. Two conventions of the gold tables are used on purpose as
target-schema knowledge: the closed vocabulary \{Frontcourt, Backcourt\} for a player's position,
and the date format \code{YYYY/M/D}.

\paragraph{Matching against gold.} Each gold row is matched to graph nodes by the article's title
and the gold name, after normalization (case, accents, punctuation), and by aliases: for people,
the same surname with a compatible first name (prefix or common nickname); for teams, the nickname
alone; for cities, \emph{City, State} and \emph{City of X}. Nodes of typed graphs (\wk{} and the
schema-based extractors) must also have the right type (\code{e1\_analysis.py}).

\section{E1: graph structure}

\paragraph{Systems.}
\begin{itemize}
  \item \textbf{\wk}, full schema, three runs.
  \item \textbf{LightRAG} 1.5.7, default extraction, chunks of 1200 tokens with 100 overlap,
    16 concurrent LLM calls (\code{lightrag\_build.py}).
  \item \textbf{GraphRAG} 3.2.0, standard indexing with its default entity types
    (organization, person, geo, event), one gleaning, chunks of 1200 tokens with 100 overlap,
    default prompts (\code{graphrag-config/}).
\end{itemize}
LightRAG and GraphRAG receive no schema.

\paragraph{Build cost.} Table~\ref{tab:cost}. Runs were executed concurrently with other runs, so
wall times are indicative.

\begin{table}[h]
\centering
\begin{tabular}{lrrrr}
\toprule
 & LLM calls & Input tokens & Output tokens & Wall time \\
\midrule
\wk{} & 4,443 & 6.5M & 1.1M & 18 min \\
LightRAG & 6,341 & 17.3M & 7.8M & 92 min \\
GraphRAG & 24,965 & 35.3M & 15.1M & 101 min \\
\bottomrule
\end{tabular}
\caption{Cost of building each graph over the 216 documents. GraphRAG's calls are
3,362 graph extractions, 16,511 description summaries and 5,092 community reports; it also used
8.2M embedding tokens.}
\label{tab:cost}
\end{table}

\paragraph{Graph.} For \wk{} we count domain nodes and edges; the provenance nodes and edges that
link every object to its source chunk and document are left out. For LightRAG we read its
GraphML graph; for GraphRAG its entity and relationship tables.

\paragraph{Relation vocabulary, own labels.} The label of an edge is what the system stores:
the relationship type in \wk, the set of keywords in LightRAG. GraphRAG stores no relation
label, only a free-text description and a weight. For a labeling we report the number of
distinct labels, the fraction used by a single edge, the number of most frequent labels needed
to cover 50\% and 90\% of the edges, and the mean number of distinct labels per pair of node
types (over pairs with at least 100 edges).

\paragraph{Relation vocabulary, derived predicates.} To measure LightRAG and GraphRAG in the same
way, and to give GraphRAG a label at all, an LLM (same model and effort as the builds) reads the
description of every edge, together with the names of its two endpoints, and names the main
relation as a short predicate in lower snake case, 1--3 words, present tense, without dates,
numbers or names, and states which endpoint is its subject (\code{derive\_labels.py}; the prompt's
examples come from an unrelated domain). Edges are labeled independently, 100 per call: 980 calls,
5.5M input and 1.5M output tokens for 97,931 edges, all of which received a label.

\paragraph{Synonym merging.} To discount surface variation, all labels are embedded
(\code{text-\allowbreak embedding-\allowbreak 3-small}, 512 dimensions) and merged by leader
clustering: in decreasing order of frequency, each label joins the most similar existing leader if
the cosine similarity is at least $t$, and otherwise becomes a leader (\code{relation\_vocab.py}).
We report $t = 0.85$. At $t = 0.80$ the largest cluster already merges opposite relations
(\code{plays\_for} with \code{plays\_against}), so 0.85 is the most aggressive threshold that still
merges only synonyms; lower thresholds are in \code{relation\_vocab.json}.

\paragraph{Direction.} For a derived predicate, an edge is stored against its relation when the
subject of the predicate is the edge's target. We also report the fraction of predicates with at
least 10 edges where at least 10\% of the edges go each way. In \wk{} direction is fixed by the
endpoint types of each relationship type. LightRAG stores an undirected graph.

\paragraph{Properties.} Typed property values per node, and the fraction of nodes whose content is
only free text (a description).

\paragraph{Control: recall against gold.} A gold attribute value (non-empty, and non-zero for
counts) is found in \wk{} if the typed field holds it (numbers within 1\%, city areas in either
square miles or km$^2$, since the gold table mixes them). For LightRAG and GraphRAG it is found if
it appears anywhere in the descriptions of the matched nodes or of their edges: dates in the usual
written forms, counts only within 60 characters of a word naming what is counted, positions through
\emph{forward/center} and \emph{guard}. This is an upper bound: the value is in the graph but not
in a form a query can select. A gold relation (player's team, team's city, team's owners, owner's
team) is found in \wk{} if an edge of the right type joins the matched nodes, and in LightRAG and
GraphRAG if any edge joins them.

\begin{table}[h]
\centering
\small
\begin{tabular}{llrrr}
\toprule
 & & \wk & LightRAG & GraphRAG \\
\midrule
Size & Nodes & 2,129 & 34,162 & 37,917 \\
 & Edges & 3,339 & 44,040 & 53,891 \\
\midrule
Node vocabulary & Node types & 4 & 14 & 99 \\
 & Node types used once & 0 & 1 & 94 \\
\midrule
Relation vocabulary, & Labels & 3 & 25,151 & none \\
own labels & Labels used once & 0\% & 79\% & -- \\
 & Labels to cover 90\% of edges & 1 & 20,747 & -- \\
\midrule
Relation vocabulary, & Labels & 3 & 4,388 & 4,428 \\
derived predicates & Labels used once & 0\% & 56\% & 52\% \\
 & Labels to cover 50\% / 90\% of edges & 1 / 1 & 51 / 1,079 & 47 / 936 \\
 & Labels per node-type pair (mean) & 1 & 310 & 736 \\
\midrule
Synonyms merged & Labels & 3 & 3,488 & 3,477 \\
($\cos \ge 0.85$) & Labels to cover 90\% of edges & 1 & 738 & 633 \\
\midrule
Direction & Edges stored against their relation & 0\% & 48\%$^{a}$ & 16\% \\
 & Frequent labels used in both directions & 0\% & 85\% & 56\% \\
\midrule
Properties & Typed property values per node & 1.5 & 0 & 0 \\
 & Nodes described only by free text & 0\% & 100\% & 100\% \\
\midrule
Control: & Entities found & 213/216 & 214/216 & 214/216 \\
recall vs.\ gold & Attribute values found$^{b}$ & 865/955 & 714/955 & 729/955 \\
 & Relations found$^{c}$ & 169/220 & 195/220 & 192/220 \\
\bottomrule
\end{tabular}
\caption{E1: structure of the three graphs. \wk{} values are means of three runs, which range
2,114--2,137 nodes, 3,278--3,390 edges, 213--214 entities, 864--867 attribute values and 168--170
relations. For \wk{} the derived and merged rows repeat its own labels. $^{a}$LightRAG's graph is
undirected. $^{b}$For LightRAG and GraphRAG, a value counts if it appears in their text (upper
bound). $^{c}$For LightRAG and GraphRAG, any edge between the two entities counts, whatever
relation it describes; most of their extra relations are players' teams outside the NBA, which the
\wk{} schema does not target.}
\label{tab:e1}
\end{table}

\paragraph{Notes.}
\begin{itemize}
  \item The three graphs find the same entities and comparable numbers of facts. \wk{} finds more
    attribute values as typed fields than the baselines hold anywhere in their text; the baselines
    connect more player--team pairs, mostly with non-NBA clubs.
  \item In the derived labelings, about 50 predicates cover half of the edges and about 1,000 are
    needed for 90\%; more than half of all predicates label a single edge. Merging synonyms at
    $t = 0.85$ reduces the 90\% figure to 738 and 633.
  \item Example: in LightRAG, the 892 edges between a player and a team in the gold tables carry 676
    distinct keyword labels; \emph{played for} labels 17 of them.
\end{itemize}

\section{X6: coupling ablation}

\paragraph{Variants.} Every key of the full schema belongs to exactly one family; variant
s$k$ keeps the core and the first $k$ families (\code{make\_coupling\_variants.py}, which checks that
s4 equals the full schema):
\begin{itemize}
  \item \textbf{s0, core:} types and their descriptions, fields (data type, description), the primary
    key of each entity type (the engine requires one), relationship endpoints, the corpus language.
    No source is bound: every entity and relationship type is extracted from every chunk of every
    collection, never at document level, and no relationship identity is declared.
  \item \textbf{s1, +bindings:} which collection and context level each type and field is found at,
    and the context levels of each endpoint.
  \item \textbf{s2, +guidance:} the corpus domain statement, instructions and examples.
  \item \textbf{s3, +validation:} options, regex, required fields.
  \item \textbf{s4, +identity (full schema):} relationship deduplication policies and merge strategies.
  \item \textbf{s5, +current\_team:} a revision of the schema made after seeing the s4 results, not an
    annotation family. In s4 a player's current team is the \code{tenure} of PlaysFor, decided by
    every chunk that mentions the player and a team; a chunk sees only part of a career, so players
    get on average about two teams marked \emph{latest}. A player's article \emph{is} the player, so
    s5 adds a Player field \code{current\_team}, read once at document level from the player's own
    article and skipped in chunks. Everything else is as in s4.
\end{itemize}

\paragraph{Measures.} On the 216 subject entities: entities found, and entities split into several
nodes. A gold attribute value is \emph{exact} if the typed field holds it in the gold's form, so a
query can compare and filter it as it is, and \emph{lenient} if it holds it in any form
(``December 30, 1984'' for 1984/12/30, ``third pick'' for 3, ``power forward'' for Frontcourt); a
value is \emph{wrong} if the field is filled but not leniently equal. Relations are the 148 gold
relations between gold entities, found when an edge of the right type joins the matched nodes. For
the current team of the 90 players whose gold team is an NBA team we report how many are found, and
the precision of what is marked as current: every team on a PlaysFor edge with tenure \emph{latest},
or the value of \code{current\_team} in s5 (\code{x6\_analysis.py}).

\begin{table}[h]
\centering
\footnotesize
\setlength{\tabcolsep}{4pt}
\begin{tabular}{lrrrrrr}
\toprule
 & s0 core & s1 & s2 & s3 & s4 full & s5 \\
 & & +bindings & +guidance & +validation & +identity & +current\_team \\
\midrule
LLM calls & 4,011 & 4,352 & 4,442 & 4,444 & 4,443 & 4,447 \\
Output tokens & 2.44M & 1.02M & 1.08M & 1.10M & 1.10M & 1.11M \\
Wall time (min) & 30 & 17 & 19 & 18 & 18 & 18 \\
Edges & 8,919 & 6,809 & 5,905 & 5,149 & 3,339 & 3,212 \\
\midrule
Entities found (of 216) & 205 & 208 & 213 & 214 & 213 & 214 \\
Entities split & 62 & 47 & 9 & 8 & 8 & 8 \\
Attribute values, exact (of 955) & 449 & 467 & 861 & 870 & 865 & 871 \\
Attribute values, lenient & 763 & 801 & 864 & 870 & 865 & 871 \\
Wrong values & 47 & 24 & 22 & 19 & 21 & 21 \\
Relations found (of 148) & 133 & 136 & 138 & 136 & 135 & 136 \\
\midrule
Current team, edges: found (of 90) & 18 & 12 & 77 & 78 & 78 & 78 \\
\quad precision & 85\% & 75\% & 53\% & 46\% & 46\% & 58\% \\
Current team, field: found (of 90) & -- & -- & -- & -- & -- & 77 \\
\quad precision & -- & -- & -- & -- & -- & 90\% \\
\bottomrule
\end{tabular}
\caption{X6: \wk{} with each schema variant, means of three runs. Over the three runs of a variant,
the values vary by at most 10 split entities, 15 attribute values, 3 relations, 2 current teams
found and 174 edges, and the current-team precision by at most 6 points, except 18 points in s0,
which marks few teams; the ranges are in \code{runs/x6/summary.csv}. Wall time excludes one run per variant from s2 to s5, during which the
computer slept.}
\label{tab:x6}
\end{table}

\paragraph{Notes.}
\begin{itemize}
  \item Without guidance, most attribute values are found but not in a form a query can use: 763
    lenient against 449 exact in s0. Guidance closes the gap (861 exact against 864 lenient in s2)
    and reduces split entities from 47 to 9.
  \item Bindings halve the output tokens (2.44M to 1.02M), shorten the build from 30 to 17 minutes and
    halve the wrong values.
  \item Identity removes repeated edges (5,149 to 3,339) without changing the scores.
  \item The current team marked on edges has low precision once guidance asks for it (46--53\%).
    Reading it as a document-level field of the player keeps the same recall and raises precision
    to 90\%.
\end{itemize}

\section{Schema-based extractors}

\paragraph{Translation rule.} Each tool receives our schema in its own format: every key of the full
schema goes into the tool's matching slot if it has one, even if the tool then ignores it, and never
into a different slot. Table~\ref{tab:slots} summarizes the mapping; each tool's \code{mapping.md}
lists every key. Each tool keeps its default chunking and pipeline; a small wrapper around its LLM
client sets the model and reasoning effort and records token usage. All three completed every
document without a failed call (\code{baselines/}).

\begin{table}[h]
\centering
\small
\begin{tabular}{llll}
\toprule
Schema element & Neo4j & LlamaIndex & OntoGPT \\
\midrule
Types, endpoints & types, patterns & names, triples & classes, slots \\
Type descriptions & yes & no slot & yes$^{*}$ \\
Field descriptions & yes & flat list over types & yes \\
Options & no slot & no slot & enum \\
Regex & no slot & no slot & pattern$^{\dagger}$ \\
Instructions & no slot & no slot & prompt annotation$^{*}$ \\
Examples & no slot & no slot & examples$^{*}$ \\
Primary key, required & KEY, EXISTENCE & entity name & key, required \\
Bindings & no slot & no slot & a root class per article type \\
Identity & exact name per type & exact name & identical id \\
Domain statement & no slot & no slot & schema description$^{*}$ \\
\midrule
Unit of extraction & 4,000 characters & 1,024 tokens & whole document \\
\bottomrule
\end{tabular}
\caption{Where each element of the schema goes. $^{*}$In the slot, but OntoGPT 1.2.0 does not put
it in the prompt. $^{\dagger}$In the template but not enforced: enforced, a single non-matching
value rejects the whole document, and 0 of 8 pilot documents survived.}
\label{tab:slots}
\end{table}

\paragraph{Neo4j \code{neo4j-graphrag} 1.22.0.} The components of \code{SimpleKGPipeline} without the
database writer: fixed-size splitter (4,000 characters, 200 overlap), LLM extractor with structured
output, schema pruning, and its default entity resolution (per type, exact match on \code{name},
first property values win, duplicate relationships merged), reimplemented in Python because it
only runs inside Neo4j with APOC. The chunk embedder is skipped; it does not affect extraction.

\paragraph{LlamaIndex \code{SchemaLLMPathExtractor} 0.14.25.} With this model and effort the stock
extractor cannot fill properties: its open property dictionary is either rejected by strict
structured output or forced to be empty. We use its documented \code{kg\_schema\_cls} parameter
with the stock models, properties typed as optional strings. Defaults otherwise: sentence splitter
(1,024 tokens, 20 overlap), at most 10 triplets per chunk, strict schema, nodes identified by name
alone.

\paragraph{OntoGPT SPIRES 1.2.0.} A LinkML template generated from the schema, one root class per
article type, selected per collection. Whole document per call, default recursion into list items,
no grounding annotators. OntoGPT splits lists on semicolons; following its own templates, the
description of each list slot starts with ``semicolon-separated list of''. Each list item is
extracted in a separate call that sees only the item, so the tenure of nearly every PlaysFor edge
is \emph{latest}.

\begin{table}[h]
\centering
\small
\setlength{\tabcolsep}{5pt}
\begin{tabular}{lrrrrr}
\toprule
 & \wk{} s0 core & \wk{} s5 & Neo4j & LlamaIndex & OntoGPT \\
\midrule
LLM calls & 4,011 & 4,447 & 2,044 & 2,295 & 2,728 \\
Output tokens & 2.44M & 1.11M & 1.37M & 1.60M & 0.39M \\
Wall time (min) & 30 & 18 & 27 & 17 & 12 \\
Edges & 8,919 & 3,212 & 4,784 & 3,590 & 2,544 \\
\midrule
Entities found (of 216) & 205 & 214 & 211 & 212 & 213 \\
Entities split & 62 & 8 & 40 & 23 & 60 \\
Attribute values, exact (of 955) & 449 & 871 & 589 & 242 & 618 \\
Attribute values, lenient & 763 & 871 & 817 & 340 & 770 \\
Wrong values & 47 & 21 & 21 & 11 & 56 \\
Relations found (of 148) & 133 & 136 & 141 & 142 & 139 \\
\midrule
Current team found (of 90) & 18 & 77 & 45 & 1 & 82 \\
\quad precision & 85\% & 90\% & 66\% & 100\% & 27\% \\
\bottomrule
\end{tabular}
\caption{Schema-based extractors, one run each, next to \wk{} with the schema core only and with
the revised full schema (means of three runs). For \wk{} s5 the current team is the document-level
field; for the others, the teams on PlaysFor edges marked \emph{latest}.}
\label{tab:baselines}
\end{table}

\paragraph{Notes.}
\begin{itemize}
  \item All three tools find the entities, and leniently most attribute values (Neo4j 817, OntoGPT 770),
    but fewer in the gold's form (589, 618), and they split 23--60 entities into several nodes.
    LlamaIndex fills few properties (340 lenient), and its name-only node identity merges entities
    of different types: it has 722 Team nodes, against 85 for \wk{} s5.
  \item OntoGPT finds 82 current teams because it marks 99\% of its PlaysFor edges as \emph{latest}.
\end{itemize}

\section{Limitations}

LightRAG, GraphRAG and the three schema-based extractors were run once; \wk{} configurations three
times. Derived predicates come from labeling each edge independently, without a shared vocabulary,
which is the situation of a user who wants to query these graphs by relation. s5 was designed
after seeing the s4 results. The gold tables contain errors (mixed units for city areas, legal names
for owners), which affect all systems alike.

\section*{Reproduction}

From \code{paper/uda-nba/}: \code{run\_full.sh} builds the LightRAG and GraphRAG graphs,
\code{run\_x6.sh} the \wk{} variants and \code{baselines/run\_baselines.sh} the schema-based
extractors. Then \code{e1\_analysis.py}, \code{derive\_labels.py} (once per RAG system),
\code{relation\_vocab.py} and \code{e1\_table.py} give Table~\ref{tab:e1}, and
\code{x6\_analysis.py} gives Tables~\ref{tab:x6} and~\ref{tab:baselines}. Data, environments and
runs are not versioned.

\end{document}
